% 图6(b) 车道汇入·MIKU：按到达时间 τ(s) 摆放同一汇入车 ST 平行四边形，
% 其左沿右移至真实进入时刻，左下方让出三角通行区，自车得以争取一段前进。
\begin{tikzpicture}[scale=1.0, transform shape, arr/.style={-Stealth, thick}]
  \useasboundingbox (-0.75,-1.05) rectangle (5.85,4.75);

  % τ(s) 左下沿（楔形）：t≥τ(s) 可达
  \draw[deepblue, very thick, densely dashed] (0,0) -- (2.6,4.2);
  \node[font=\tiny, gray!55!black, align=center] at (0.55,3.75) {$t{<}\tau(s)$\\不可达};

  % 汇入车 ST 轨迹：同一平行四边形，左沿右移至 t=1.5（真实进入时刻）
  \fill[dangerred!35, draw=dangerred, thick]
    (1.5,1.2) -- (5.0,3.0) -- (5.0,4.0) -- (1.5,2.2) -- cycle;
  \node[font=\scriptsize\bfseries, dangerred!85!black, rotate=22] at (3.3,2.9) {汇入车 ST 轨迹};
  \draw[<-, dangerred!75!black, thin] (0.4,1.7) -- (1.45,1.7);
  \node[font=\tiny, dangerred!80!black, anchor=east, align=right] at (0.4,1.7) {左沿\\按 $\tau$\\右移};

  % 争取出的三角通行区（左下方）
  \fill[lightgreen!35] (0,0) -- (1.5,1.2) -- (1.5,0) -- cycle;
  \node[font=\tiny\bfseries, lightgreen!45!black, align=center] at (0.95,0.45) {通行区};

  % 坐标轴
  \draw[arr] (0,0) -- (5.6,0) node[right, font=\scriptsize] {$t$/s};
  \draw[arr] (0,0) -- (0,4.5) node[above, font=\scriptsize] {$s$/m};
  \node[font=\tiny, left=2pt] at (0,0) {$0$};

  % MIKU s(t)：沿通行区争取前进，随后与平行四边形下沿平行前进但保持间隙、不贴合
  \draw[deeppink, very thick, -Stealth]
    (0,0) .. controls (0.7,0.7) and (1.1,0.85) .. (1.4,0.9)
    .. controls (2.6,1.35) and (3.8,1.95) .. (5.0,2.45);
  \fill[deeppink] (1.4,0.9) circle (1.6pt);
  \node[font=\tiny, deeppink, anchor=north west] at (1.55,0.85) {争取前进后平行跟行};
  % 与下沿的间隙指示
  \draw[<->, gray!60, thin] (3.5,1.72) -- (3.5,2.28);
  \node[font=\tiny, gray!55!black, anchor=west] at (3.55,2.0) {间隙};
\end{tikzpicture}
